% Paper Summary: LLaMA (Touvron et al., 2023)

\begin{papersummary}{2023}{LLaMA: Open and Efficient Foundation Language Models}{Hugo Touvron, Thibaut Lavril, Gautier Izacard, Xavier Martinet, Marie-Anne Lachaux, Timothée Lacroix, Baptiste Rozière, Naman Goyal, Eric Hambro, Faisal Azhar, Aurelien Rodriguez, Armand Joulin, Edouard Grave, Guillaume Lample}{LLaMA demonstrated that smaller models trained far beyond the compute-optimal token budget match much larger models at a fraction of the inference cost, and its open release catalyzed open-source LLM development.}

\summaryheading{Key Ideas}
LLaMA extended the Chinchilla analysis to account for inference cost: because serving cost depends on model size, a model destined for deployment should be trained past the compute-optimal point. Its models, from 7B to 65B parameters trained on 1--1.4T tokens---well beyond the compute-optimal budget at the smaller sizes---achieved performance competitive with GPT-3 (175B) and Chinchilla (70B). The models used only publicly available data, demonstrating that frontier performance doesn't require proprietary datasets. LLaMA employed standard Transformer architecture with modern improvements (\hyperref[paper:su-2021]{RoPE}, \hyperref[paper:shazeer-2020]{SwiGLU}, \hyperref[paper:zhang-2019-rmsnorm]{RMSNorm}) showing that careful training matters more than exotic architectures. The open release enabled research and applications previously limited to organizations with massive resources.

\summaryheading{Follow-on Works}
Llama 2 and Llama 3 continued the series with improved training and safety. Alpaca, Vicuna, and countless fine-tuned variants built on LLaMA. The model catalyzed open-source LLM development including Mistral, Falcon, and others. LLaMA's architecture and training approach influenced numerous subsequent models. The democratization of frontier-class models enabled widespread experimentation and deployment.

\summaryheading{Lasting Contributions}
LLaMA democratized access to frontier-class language models, enabling organizations and researchers without massive compute budgets to work with capable models. The open release catalyzed an explosion of research, fine-tuning, and applications. LLaMA demonstrated that the ``secret sauce'' of frontier models is primarily careful, extended training on quality data rather than proprietary techniques. The model's influence on the open-source ecosystem makes it among the most impactful LLM releases, shifting power dynamics in AI development and enabling the diverse LLM applications we see today.

\end{papersummary}
